% !TeX root = ../main.tex

\chapter{相关背景}

\section{大语言模型与代码任务}

\subsection{大语言模型的能力演进}

大语言模型是以 Transformer~\cite{vaswani2017attention} 为基础、通过大规模语料预训练得到的神经网络模型，具备自然语言理解、生成与一定推理能力。自 GPT-3~\cite{brown2020language} 展示出显著的上下文学习能力以来，大语言模型在问答、摘要、翻译、推理和代码生成等任务中均取得了快速发展。随后，GPT-4~\cite{achiam2023gpt}、Gemini~\cite{team2023gemini}、Llama 2~\cite{touvron2023llama} 等新一代模型进一步增强了复杂推理、长上下文处理与多工具交互能力，推动了大模型从传统自然语言处理逐步进入代码、数学和多模态任务场景。

从软件工程角度看，大语言模型的意义不仅在于“能够生成代码”，更在于其具备跨模态、跨阶段的问题求解潜力。它们既能够理解自然语言描述的问题，又能够处理结构化的程序文本，还可以在一定程度上结合外部工具完成更复杂的分析与决策。这为自动程序修复、代码理解与软件工程智能体的发展奠定了基础。

\subsection{大语言模型在代码任务中的应用}

代码是大语言模型最重要的应用领域之一。Codex~\cite{chen2021evaluating} 的出现表明，通过在大规模开源代码上进行训练，大语言模型能够较好地完成代码补全、函数合成与 API 调用预测等任务。随后，CodeBERT~\cite{feng2020codebert}、CodeT5~\cite{wang2021codet5}、DeepSeek-Coder~\cite{guo2024deepseekcoder} 等模型进一步提升了代码理解与生成能力，使得大模型在代码补全、代码翻译、代码摘要、漏洞定位与程序修复等任务中表现出越来越强的实用价值。

值得注意的是，代码任务本身也在发生演进。早期研究主要关注函数级或片段级任务，而随着实际需求提升，研究对象逐渐扩展到跨文件代码补全、仓库级代码问答、GitHub issue resolution 等更接近真实软件工程的问题形态。在这些任务中，模型不再只需要补全局部代码，而需要在更大范围内理解程序结构、项目上下文和任务描述。

\subsection{大语言模型在复杂代码任务中的局限性}

尽管大语言模型在代码任务中取得了显著进展，但其在复杂软件工程场景中仍面临若干关键限制。首先，\textbf{上下文容量仍是约束因素}。真实代码仓库往往由大量文件组成，即便模型支持较长上下文窗口，也很难在单次输入中完整覆盖所有相关信息。其次，\textbf{代码场景中的幻觉问题依然存在}，模型可能生成语法正确但语义错误的代码，或错误引用不存在的 API 与标识符。再次，\textbf{多步任务的求解成本迅速增长}。在需要跨文件理解、环境执行与多轮决策的任务中，模型需要不断读取新信息并维护历史轨迹，这使复杂代码任务的开销远高于单轮生成场景。

这些限制表明，大语言模型在简单代码生成任务中的成功并不能直接推广到真实仓库级软件工程任务。要让大模型在复杂场景中真正可用，需要进一步引入工具调用、环境交互和上下文管理机制。


\section{代码智能体的基本工作机制}

\subsection{代码智能体的定义与特点}

代码智能体（Code Agent）是以大语言模型为推理核心，并具备工具调用与环境交互能力的自动化软件系统。与传统单轮代码生成模型相比，代码智能体最重要的特征在于其\textbf{过程性}与\textbf{交互性}：它不是一次性输出最终代码结果，而是在多轮感知—决策—行动循环中逐步推进任务求解。

一般而言，代码智能体具备若干常见操作能力，包括：在仓库中搜索相关文件或函数、读取指定文件内容、编辑目标代码、执行命令或测试脚本、分析运行反馈并据此调整后续策略。这些能力使其能够处理单步生成模型难以解决的复杂任务，例如真实 issue 修复、代码审查、环境调试与测试驱动修改等。

\subsection{多轮交互决策与过程性成本}

代码智能体的执行过程本质上是一个连续的决策序列。在每一轮中，智能体根据当前任务描述、历史执行轨迹以及最新工具返回结果，决定下一步是继续搜索、读取文件、修改代码，还是运行测试进行验证。执行结束后，本轮的 observation 又会被纳入后续轮次的上下文中，成为下一次推理的输入。

这种机制显著增强了智能体的问题求解能力，但同时也引入了额外成本。随着轮次增加，历史文件内容、命令输出、错误信息和中间推理不断累积，后续模型调用的输入规模也随之膨胀。尤其在仓库级任务中，代码搜索和文件阅读往往会引入大量长文本，而多次失败尝试还会带来冗余 observation 与重复操作。因此，代码智能体的成本不应只看最后一次模型调用，而应从全过程角度理解。

\subsection{典型代码智能体框架}

近年来，围绕代码智能体出现了多种具有代表性的框架。SWE-agent~\cite{yang2024sweagent} 将自然语言任务描述、工具接口与环境操作相结合，展示了 agent-computer interface 在软件工程自动化中的潜力；AutoCodeRover~\cite{zhang2024autocoderover} 强调通过更结构化的仓库搜索与定位机制提升自动程序修复能力；Agentless~\cite{xia2024agentless} 则从相反方向切入，分析复杂 agent 框架本身带来的额外开销，探索更轻量的代码问题求解路径。这些工作共同说明，代码智能体的研究不仅关注模型能力，也 increasingly 关注搜索策略、上下文组织方式以及系统开销控制。


\section{自动程序修复与仓库级修复任务}

\subsection{自动程序修复的基本目标}

自动程序修复（Automated Program Repair，APR）的目标是在尽可能少的人为干预下，自动生成能够修复软件缺陷的补丁。传统 APR 方法主要包括基于搜索的修复~\cite{le2011genprog}、基于语义约束的修复~\cite{nguyen2013semfix} 和基于模板的修复~\cite{liu2019tbar} 等。它们通常依赖预定义的变异空间、补丁模板或约束求解技术，在一定范围内取得了成功，但在复杂、开放世界的软件项目中仍存在适用范围受限的问题。

大语言模型的出现为 APR 带来了新的可能性。与传统方法相比，基于大语言模型的修复方法能够更自然地结合自然语言问题描述、代码上下文与测试反馈，从而直接生成候选补丁或辅助完成修复流程。随着模型能力增强，APR 的研究重点也逐渐从“如何在局部代码上生成补丁”扩展到“如何在真实仓库环境中完成端到端问题修复”。

\subsection{仓库级修复任务的典型流程}

在仓库级场景下，自动程序修复通常包括以下几个关键阶段：（1）\textbf{问题理解}，即根据 issue 描述、测试失败信息或报错日志理解待修复问题；（2）\textbf{代码定位}，即在项目仓库中搜索与问题相关的文件、函数或调用链；（3）\textbf{上下文组织}，即选择并整理与修复相关的代码片段和环境信息；（4）\textbf{补丁生成}，即调用模型生成候选修改；（5）\textbf{验证与筛选}，即通过测试运行、静态检查或人工规则筛除无效补丁。

这一流程与传统单文件代码生成存在显著差异。单文件生成任务通常将问题与全部上下文一次性提供给模型，而仓库级修复任务需要系统在多轮交互中主动发现、构建并维护上下文，因此其任务链路更长，对上下文选择和过程管理的依赖更强。

\subsection{仓库级修复任务的特殊挑战}

仓库级修复任务的主要挑战可概括为以下几个方面。

\textbf{第一，搜索空间庞大。} 真实软件仓库可能包含数百甚至上千个文件，问题涉及的位置并不明确，智能体必须在有限代价下尽快缩小搜索范围。

\textbf{第二，跨文件依赖复杂。} 一个 bug 往往并不局限于单个函数或单个文件，而是与调用关系、继承结构、配置逻辑或测试脚本共同相关。理解一处问题，常常需要追溯多个文件的语义联系。

\textbf{第三，上下文组织难度高。} 在有限上下文窗口下，哪些文件该读、哪些代码该保留、哪些 observation 应当被压缩或丢弃，直接决定了模型后续推理效果。

\textbf{第四，验证成本显著。} 仓库级修复通常需要真实环境执行测试，而测试反馈又会反过来进入后续上下文，导致整体任务开销进一步放大。

这些特点决定了仓库级修复任务不仅是一个补丁生成问题，更是一个复杂的过程管理与成本控制问题。


\section{GitHub Issue Resolution 与 SWE-bench 场景}

\subsection{从代码生成到 GitHub Issue Resolution}

随着大语言模型在代码任务中的发展，研究对象逐渐从函数级生成扩展到更真实的软件工程问题。其中，GitHub issue resolution 是一类具有代表性的任务：系统需要根据 issue 描述理解问题，在给定仓库中定位相关代码，并最终生成可执行、可验证的解决方案。与传统代码补全相比，这类任务不仅涉及代码生成，还需要处理自然语言需求理解、跨文件上下文选择、环境交互与结果验证等环节。

围绕这一任务，研究社区提出了多种 benchmark。除了聚焦代码仓库理解的基准外，OmniGIRL~\cite{guo2025omnigirlmultilingualmultimodalbenchmark} 进一步从多语言、多模态 GitHub issue resolution 的角度扩展了任务边界，说明真实问题求解并不局限于单一语言或纯文本输入。这些工作共同表明，GitHub issue resolution 已成为评估大模型软件工程能力的重要方向。

\subsection{SWE-bench 的任务形式与评测意义}

在现有问题解决基准中，SWE-bench~\cite{jimenez2024swebench} 是最具代表性的仓库级代码修复 benchmark 之一。该基准从多个真实开源 Python 项目中收集 issue、仓库快照与验证测试，将“是否能够在给定仓库中生成正确补丁并通过测试”作为核心评测目标。与函数级或单文件任务相比，SWE-bench 更强调真实项目中的搜索、理解、修改与验证链路，因此更接近代码修复智能体在工程实践中的使用方式。

此后，研究者又提出了 SWE-bench Verified 这一经人工筛选和校验的高质量子集，以减少原始 benchmark 中的噪声与标注不确定性，提升评测的稳定性。这使得 SWE-bench 及其变体逐渐成为研究仓库级自动程序修复和代码智能体的重要标准场景。

\subsection{SWE-bench 适合研究推理成本的原因}

对于本文而言，SWE-bench 的价值不仅在于可评估修复成功率，更在于它非常适合用于研究代码修复智能体的推理成本。首先，该 benchmark 的任务通常需要多轮搜索、代码阅读和环境验证，使得全过程成本具有充分的分析空间。其次，不同任务在搜索范围、代码依赖复杂度和验证反馈长度上差异明显，便于观察成本在不同样本上的分布规律。再次，SWE-bench 的输出结果既可以用“是否修复成功”衡量，又可以与 token 开销、运行时延等成本指标联合分析，使其天然适合研究能力与成本之间的权衡关系。

因此，尽管 OmniGIRL 等工作从更广义的 issue resolution 角度拓展了任务生态，本文仍以 SWE-bench 作为核心实验场景，将研究焦点集中于仓库级代码修复智能体的高推理成本问题。


\section{Token、上下文窗口与推理成本}

\subsection{Token 的基本概念}

Token 是大语言模型处理文本的基本单位，通常对应若干字符构成的子词片段。在代码场景中，标识符、关键字、运算符和自然语言注释都会被分解为不同 token。对于 API 模型而言，调用成本通常按输入 token 与输出 token 分别计费，因此 token 数量可以看作推理成本最直接、最细粒度的计量单位。

对于代码修复智能体而言，输入 token 的来源包括系统提示、任务描述、代码上下文、工具返回结果与历史交互轨迹；输出 token 则包括模型生成的中间推理、工具调用指令和补丁内容。在多轮交互任务中，输入 token 的累计开销通常远大于输出 token。

\subsection{长上下文的成本效应}

上下文窗口决定了模型单次调用能够接收的最大 token 数量。尽管现代模型的上下文长度较早期系统已有大幅提升，但更长的上下文并不意味着成本问题自动消失。相反，在长上下文代码任务中，更多内容能够被“装进去”，也意味着更高的调用费用与更大的延迟开销。

对于多轮交互式任务而言，长上下文的成本效应更加明显。若系统在每一轮都携带完整的历史上下文，则随着轮次增加，后续调用会承载越来越多的历史信息，使累计输入成本迅速增长。此前关于长上下文代码压缩的研究已经表明，代码场景中的压缩通常伴随性能与效率之间的显著权衡，这说明长上下文并不仅是“能否放下”的问题，更是“如何有选择地组织和保留”的问题。

\subsection{过程性推理成本的定义}

为了刻画代码智能体的真实开销，本文将推理成本理解为\textbf{过程性成本}。设一个修复任务共经历 $N$ 轮模型调用，第 $k$ 轮调用的输入 token 数为 $c_k^{in}$，输出 token 数为 $c_k^{out}$，则其总 token 成本可抽象表示为：

\[
C_{total} = \sum_{k=1}^{N}(c_k^{in} + c_k^{out})
\]

若进一步考虑不同 API 对输入与输出 token 的不同计费权重，则还可在此基础上引入价格系数进行扩展。该定义强调两点。第一，代码智能体的成本并非单次调用成本，而是整个任务链路的累计结果；第二，成本会随着搜索、阅读、验证与历史 observation 的累积而动态变化，因此分析对象应从“最终 prompt 长度”扩展到“全过程执行轨迹”。

\subsection{成本的扩展维度}

除 token 开销外，代码修复智能体的实用成本还包括运行时延、交互轮次数以及单位成功成本等维度。运行时延直接影响用户体验与系统并发能力；交互轮次数能够反映一个智能体是否存在低收益反复尝试；单位成功成本则衡量平均修复一个 issue 所需的真实代价。在工程部署场景中，这些因素往往与 token 成本共同决定一个系统是否真正可用。因此，本文后续将围绕 token 成本展开分析，同时在讨论中兼顾其与其他成本维度之间的关系。
